\chapter{Conclusions and Recommendations}

\section{Conclusions}
The Thesis and Project Management System (TPMS) was successfully designed, developed, and tested as a centralized, role-aware web platform for managing the complete academic lifecycle of Bachelor projects (Minor and Major) and Master research theses at the Department of Electronics and Computer Engineering (DOECE), Pulchowk Campus, Institute of Engineering, Tribhuvan University.

The system fulfilled all its established objectives:
\begin{enumerate}
  \item Digitalized manual, paper-based project and thesis management workflows into a single unified web platform.
  \item Provided role-specific dashboards for Students, Supervisors, External Examiners, Program Coordinators, and Maintainers with clear milestone visibility.
  \item Implemented automated group formation, supervisor/examiner allocation mechanisms, and bulk student onboarding via Excel spreadsheets.
  \item Automated criteria-based defense evaluation across departmental scoring rubrics, eliminating manual arithmetic errors and generating official print-ready A4 PDF evaluation sheets.
  \item Integrated automated notification mechanisms to keep all academic stakeholders informed of key milestones, allocations, and defense updates.
\end{enumerate}

Through systematic functional and integration testing, the platform proved capable of streamlining administrative tasks, enforcing academic integrity rules, and reducing the operational workload required for academic project governance.

\section{Recommendations and Future Enhancements}
Based on the implementation experience and departmental feedback, the following future enhancements are recommended:

\begin{enumerate}
  \item \textbf{Defense Scheduling and Calendar Integration:} Implement an automated defense scheduling module that allows coordinators to configure defense time slots, enables evaluators to confirm availability, and generates official defense timetables.
  
  \item \textbf{Configurable Evaluation Rubric Builder:} Develop an administrative interface allowing coordinators to dynamically customize rubric criteria, weightages, and score boundaries for different academic sessions and project types.
  
  \item \textbf{Progressive Web App (PWA) and Mobile Support:} Package the web application as a Progressive Web App (PWA) with responsive mobile optimization and push notifications to enhance accessibility for students and faculty members.
  
  \item \textbf{Multi-Institution Tenancy:} Re-architect the database schema and API layer to support multi-tenant deployment, enabling other IOE campus departments to independently administer their project management workflows within a shared platform instance.
\end{enumerate}
